
\illustration{0.4}{harddrive}{Image taken from \href{https://commons.wikimedia.org/wiki/File:ISS_laptop_hard_drive_failure_error_message.jpg}{commons.wikimedia.org}.}

In an earlier contest, contestants had to solve a simple
problem. They were given a string and had to print each
unique substring of it, along with the number of occurrences
it had in the original string. For example \texttt{AB}
would print \texttt{A 1 B 1 AB 1} and \texttt{AAA} would
print \texttt{A 3 AA 2 AAA 1}.

When copying over this problem for reuse in this contest,
several mistakes were made. The input constraints were changed
significantly, making the problem absolutely impossible!
Luckily this was partially cancelled out by the output
validator being mangled as well. Now instead of checking
for absolute correctness it only requires the number of
occurrences of each character in the output to be correct.
With a quick hotfix of applying run length encoding to the
output, the problem was finally solvable again and the
contest could continue on as planned. Right?

